\documentclass[11pt,a4paper]{article}
\usepackage{iftex}
\usepackage{amsmath,amsthm}
\ifPDFTeX
  \usepackage[T1]{fontenc}
  \usepackage[utf8]{inputenc}
  \usepackage{textcomp}
  \usepackage{amssymb}
\else
  % XeTeX/LuaTeX: amssymb BEFORE unicode-math (unicode-math then overrides it cleanly;
  % the reverse order clashes on \eth). Keeps \square, \mathbb, \mathfrak available.
  \usepackage{amssymb}
  \usepackage{unicode-math}
  \defaultfontfeatures{Scale=MatchLowercase}
\fi
\usepackage{booktabs}
\usepackage{longtable}
\usepackage{array}
\usepackage{graphicx}
\usepackage[margin=1in]{geometry}
\usepackage{xcolor}
\usepackage{hyperref}
\hypersetup{colorlinks=true,linkcolor=blue!60!black,citecolor=green!50!black,urlcolor=blue!60!black}
\usepackage{fancyhdr}
\pagestyle{fancy}
\fancyhf{}
\fancyhead[L]{\small PACSeries Paper 10}
\fancyhead[R]{\small Dawn Field Institute}
\fancyfoot[C]{\thepage}
\renewcommand{\headrulewidth}{0.4pt}
\setlength{\parindent}{0pt}
\setlength{\parskip}{6pt plus 2pt minus 1pt}
\providecommand{\tightlist}{\setlength{\itemsep}{0pt}\setlength{\parskip}{0pt}}

\title{Connection, Identity, and Spacetime}
\author{Peter Groom \\ Dawn Field Institute}
\date{May 2026}

\begin{document}
\maketitle

\subsection*{On deriving the Lorentz group, invariant interval, and speed of light from ADE graph structure and SEC complexification}

\textbf{Peter Groom, Dawn Field Institute} \\
\textbf{PACSeries Paper 10} \\
\textbf{Date}: May 2026 \\
\textbf{Version}: 1.0 (Draft)

\begin{center}\rule{0.5\linewidth}{0.4pt}\end{center}

\section*{Abstract}

We derive the Lorentz group, the invariant interval, and the speed of light from the graph-theoretic foundations of Dawn Field Theory (DFT). The central claim is that connection is addition: the operation of connecting two nodes on a graph is algebraically identical to addition on their label sets, and the graphs that close under this operation are precisely the ADE Dynkin diagrams — a classification theorem from mathematics, not a postulate of physics.

Milestone 12 of DFT establishes three results. (1) The ADE classification is equivalent to Fibonacci arithmetic closure: the only ADE types compatible with Fibonacci recursion are $A_1$ (trivial), $D_4$ (SU(2)), and $E_8$ (containing SU(3)). The Standard Model gauge group SU(3) $\times$ SU(2) $\times$ U(1) is not assumed; it is the unique Fibonacci-compatible content of ADE. The gauge dimension closure $F_7 = 13 = 1 + 3 + 8 + 1$ follows from $\dim(\text{U}(1)) + \dim(\text{SU}(2)) + \dim(\text{SU}(3)) + 1 = F_7$. (2) Physical laws emerge as basin attractors of information dynamics, with relaxation ratios $\varphi^{d_2 - d_1}$ between force depths, converged to rather than imposed. (3) SEC complexification of the simplest ADE type ($A_1$) produces $A_1 \to \text{SL}(2,\mathbb{C}) \cong \text{SO}(3,1)$, the Lorentz group, derived from connection plus entropy.

Milestone 13 extends this to spacetime geometry. (4) Identity is complement: the identity of a vertex on an ADE graph is its complement spectrum, the list of what it is not. This is an outside-in definition, requiring no intrinsic properties. (5) Different observers compute different complements of the same target: this is definitional parallax, and complement-transformations between observers form the Weyl group of the graph. (6) SEC extends $\mathbb{Z}_2$ (graph complement) to $\text{SL}(2,\mathbb{C})$ (continuous deformation) to $\text{SO}(3,1)$ (Lorentz). (7) The invariant interval $ds^2$ emerges from the uniqueness of the Killing form on the Lie algebra, the only bilinear form invariant under all complement transformations. (8) The speed of light is the coherence limit on complement-deformation: the maximum rate at which one vertex's complement can change while maintaining graph consistency. Proper time follows: $d\tau = dt / \cosh(\eta)$ where $\eta$ is the boost parameter.

The complete derivation chain — self-loop $\to$ $\varphi$ $\to$ PAC $\to$ ADE $\to$ complement $\to$ parallax $\to$ Weyl $\to$ SEC $\to$ Lorentz $\to$ $ds^2$ $\to$ $c$ $\to$ proper time — contains 12 links, each computationally verified. We report 49/52 (94\%) on M12 and 51/60 (85\%) on M13/13.5, with 8 failures documented in full (including PSD degeneracy proven fundamental and rapidity composition with discrete gap). 16 predictions (8 from each milestone), zero contradictions with M1–M11.

\textbf{Keywords}: ADE classification, connection, Lorentz group, SEC complexification, complement identity, Weyl group, invariant interval, speed of light, spacetime derivation, Dawn Field Theory

\textbf{Epistemic status.} This paper is part of the PACSeries, an evolving research program rather than a finished theory (see the series note \emph{How to Read the PACSeries}). Its quantitative claims are tiered by derivation type — \textbf{A} structural, \textbf{B} identified, \textbf{C} pattern-matched — and labeled in the text. Failures are reported rather than omitted, and framework-level corrections are logged in the Epistemic Corrections Registry in the repository. Every headline result is reproducible from the \texttt{Code/} and \texttt{Data/} in this package (provenance in \texttt{trace.yaml}).

\begin{center}\rule{0.5\linewidth}{0.4pt}\end{center}

\section*{1. The connection problem}

Spacetime has two features that every physicist uses but no standard theory derives:

\begin{enumerate}\tightlist
\item \textbf{The Lorentz group SO(3,1)} — why this symmetry, and not some other?
\item \textbf{The invariant interval $ds^2 = -c^2 dt^2 + dx^2 + dy^2 + dz^2$} — why this metric signature?
\end{enumerate}

General relativity takes the Lorentz group as given, quantum field theory builds on it, and string theory embeds it; none derives it from anything simpler.

This paper derives both from the combinatorial structure of ADE graphs — the same classification that already governs gauge symmetries, root systems, and singularity resolutions in mathematics.

\begin{center}\rule{0.5\linewidth}{0.4pt}\end{center}

\section*{Part I: Connection as Primitive (Milestone 12)}

\section*{2. Connection is addition}

\subsection*{2.1 The algebraic identity}

Consider a graph $G$ with vertices $V$ and edges $E$. The operation of \emph{connecting} two vertices — asking whether and how they communicate — is equivalent to the algebraic operation of addition on the label space.

For ADE Dynkin diagrams, the adjacency matrix $A$ encodes which vertices connect. The graph Laplacian $L = D - A$ (degree matrix minus adjacency) captures the dynamics of information flow. The eigenvalues of $L$ are the resonant frequencies of the graph; the eigenvectors are the modes.

For ADE graphs (M12 exp\_01), the operation "connect node $i$ to node $j$" is isomorphic to addition in the root lattice. This is an algebraic identity, proven by the classification theorem.

\subsection*{2.2 ADE as Fibonacci arithmetic closure}

Which ADE types are compatible with Fibonacci recursion $F_n = F_{n-1} + F_{n-2}$?

The Fibonacci sequence modulo the Coxeter number $h$ of each ADE type determines compatibility (M12 exp\_03). A type is Fibonacci-compatible if the Fibonacci orbit under mod $h$ visits all residues — i.e., if Fibonacci arithmetic is complete on that graph.

Result: exactly three families are Fibonacci-compatible:
\begin{itemize}\tightlist
\item $A_1$ ($h = 2$): trivial, $\dim = 0$ (U(1) contribution)
\item $D_4$ ($h = 6$): $\dim = 3$, corresponds to SU(2)
\item $E_8$ ($h = 30$): $\dim = 8$, contains SU(3) as maximal subgroup
\end{itemize}

The gauge dimension closure:

\[F_7 = 13 = \underbrace{1}_{\text{U(1)}} + \underbrace{3}_{\text{SU(2)}} + \underbrace{8}_{\text{SU(3)}} + \underbrace{1}_{\text{singlet}}\]

This is not numerology. $F_7 = 13$ is the seventh Fibonacci number; the dimensions 1, 3, 8, 1 are forced by the ADE classification; their sum being exactly $F_7$ is the Fibonacci closure condition.

\subsection*{2.3 Interpretation}

The Standard Model gauge group is the unique content of ADE arithmetic compatible with Fibonacci recursion, which is itself the unique recursion compatible with PAC conservation (Paper 1).

\begin{center}\rule{0.5\linewidth}{0.4pt}\end{center}

\section*{3. Laws as basin attractors}

\subsection*{3.1 Convergence, not imposition}

Physical laws are basin attractors of information dynamics, not rules imposed on nature (M12 exp\_03–05).

Start with random initial conditions on an ADE graph. Apply PAC-conserving dynamics (total information conserved at every step). The system converges to a fixed point corresponding to a physical law, not because the law was encoded but because the dynamics have no alternative.

The relaxation ratios between force depths follow $\varphi^{d_2 - d_1}$:
\begin{itemize}\tightlist
\item Strong/EM ratio: $\varphi^{13-3} = \varphi^{10} \approx 122.99$ (vs observed $\alpha_{\text{EM}}/\alpha_s \sim 130$)
\item EM/weak ratio: $\varphi^{7-13}$ (inverse Fibonacci depth ordering)
\end{itemize}

\subsection*{3.2 Thermodynamic mechanism}

The convergence has a thermodynamic explanation (M12 exp\_06–07). Each ADE basin has an entropy landscape. The free energy $F = E - TS$ has a unique minimum at the physical law. The depth of this minimum — and therefore the stability of the law — is proportional to the cascade depth.

Gravity (depth 183) has the shallowest basin per unit coupling. This is why gravity is "weakest": not because of a hierarchy problem but because depth-183 basins are the most extended in configuration space.

\begin{center}\rule{0.5\linewidth}{0.4pt}\end{center}

\section*{4. SEC complexification: Lorentz from entropy}

\subsection*{4.1 The derivation}

The simplest ADE type is $A_1$: a single edge connecting two vertices. Its symmetry group is $\mathbb{Z}_2$ (swap the vertices).

SEC (Symbolic Entropy Collapse) complexifies this symmetry. The physical reasoning: entropy flow is directional (second law), which breaks the discrete $\mathbb{Z}_2$ into a continuous flow. The mathematical consequence:

\[A_1 \xrightarrow{\text{SEC}} \mathfrak{sl}(2, \mathbb{C}) \cong \mathfrak{so}(3,1)\]

Step by step (M12 exp\_10–11):
\begin{enumerate}\tightlist
\item $A_1$ has Lie algebra $\mathfrak{su}(2)$ (3 generators: rotations)
\item SEC promotes real parameters to complex: $\theta \to \theta + i\eta$
\item Complexification: $\mathfrak{su}(2)_{\mathbb{C}} = \mathfrak{sl}(2, \mathbb{C})$
\item Real form: $\mathfrak{sl}(2, \mathbb{C}) \cong \mathfrak{so}(3,1)$ (Lorentz algebra)
\end{enumerate}

The Lorentz group is not postulated here; it follows from the simplest graph symmetry once the SEC arrow is imposed. All 15 independent commutation relations of $\mathfrak{so}(3,1)$ close to machine precision (maximum closure error $1.1\times10^{-16}$; M12 exp\_11): $[J_i,J_j]=i\epsilon_{ijk}J_k$, $[K_i,K_j]=-i\epsilon_{ijk}J_k$, and $[J_i,K_j]=i\epsilon_{ijk}K_k$. The minus sign in the boost-boost bracket, which distinguishes Minkowski from Euclidean signature, follows from $K_i=iJ_i$ rather than being imposed.

\subsection*{4.2 Why this works}

The isomorphism $\text{SL}(2, \mathbb{C}) \cong \text{SO}(3,1)$ is a known mathematical fact. What DFT adds is the \emph{physical mechanism} for complexification: SEC dynamics require a preferred direction (actualization over potential), which breaks discrete symmetry into continuous flow, which forces complexification.

The 3 rotation generators become 3 rotations + 3 boosts = 6 generators of the Lorentz group. No dimension is added by hand; it emerges from the real/imaginary decomposition of complexified $\mathfrak{su}(2)$.

\begin{center}\rule{0.5\linewidth}{0.4pt}\end{center}

\section*{Part II: Identity as Complement (Milestone 13)}

\section*{5. Identity is what you are not}

\subsection*{5.1 The complement spectrum}

On an ADE graph, every vertex $v$ has a complement $\bar{v}$ — the set of all vertices not adjacent to $v$. The complement spectrum of $v$ is the multiset of distances from $v$ to all other vertices (M13 exp\_01).

The complement spectrum uniquely identifies the vertex. You do not need to know what a vertex is; you can reconstruct it entirely from what it is not.

This is identity as complement: an outside-in definition. It requires no intrinsic properties, no labels, no coordinates. It is purely relational.

\subsection*{5.2 Definitional parallax}

Different vertices compute different complements of the same target (M13 exp\_02). If vertex $A$ defines vertex $B$ through its complement, and vertex $C$ also defines vertex $B$ through its complement, the two definitions agree on the \emph{identity} of $B$ but disagree on the \emph{description}.

This disagreement is definitional parallax, the graph-theoretic origin of reference frames. A "reference frame" is a choice of which vertex does the complement computation.

\begin{center}\rule{0.5\linewidth}{0.4pt}\end{center}

\section*{6. From complement to Lorentz}

\subsection*{6.1 Complement-transformations form the Weyl group}

The transformations that relate different complement descriptions of the same vertex form a group — the Weyl group of the ADE graph (M13 exp\_04).

For $A_1$: Weyl group = $\mathbb{Z}_2$ (swap the two vertices) For $D_4$: Weyl group = $S_3$ (permute the three leaves) For $E_8$: Weyl group = $W(E_8)$ (order 696,729,600)

The Weyl group acts on complement descriptions the way Lorentz transformations act on spacetime coordinates: this is the discrete precursor to continuous Lorentz symmetry.

\subsection*{6.2 SEC extends the Weyl group}

SEC complexification extends the discrete Weyl group to a continuous Lie group (M13 exp\_05):

\[\mathbb{Z}_2 \xrightarrow{\text{SEC}} \text{SL}(2, \mathbb{C}) \xrightarrow{\text{real form}} \text{SO}(3,1)\]

This is the same derivation as Section 4, but now with a geometric interpretation: the Lorentz group is the continuous version of complement-transformations between observers.

\subsection*{6.3 The invariant interval}

If complement-transformations are the Lorentz group, what is invariant under them?

The Killing form of $\mathfrak{sl}(2, \mathbb{C})$ is the unique bilinear form invariant under all adjoint transformations (M13 exp\_09). Its signature is $(-,+,+,+)$. This is the spacetime metric:

\[ds^2 = -c^2 dt^2 + dx^2 + dy^2 + dz^2\]

The metric signature is not chosen; it is forced by the Killing form of the complexified $A_1$ algebra. On the 6-dimensional real algebra the Killing form has signature $(3,3)$ (rotations positive, boosts negative); restricted to the 4-dimensional vector representation it has signature $(1,3)$, the Minkowski metric (M13 exp\_09). The interval $s^2 = x^{\top}\eta\,x$ is preserved across 135 Lorentz transforms (maximum relative error $2.9\times10^{-6}$, growing with rapidity up to $\eta=10$), and the invariant form is unique: the space of symmetric bilinear forms preserved by all $\mathfrak{so}(3,1)$ generators is one-dimensional (Schur's lemma), so the metric is a consequence of the algebra rather than a choice.

\begin{center}\rule{0.5\linewidth}{0.4pt}\end{center}

\section*{7. Speed of light and proper time}

\subsection*{7.1 The coherence limit}

The speed of light is the maximum rate at which complement-deformation can propagate while maintaining graph consistency (M13 exp\_07–08).

On an ADE graph, changing the complement of vertex $v$ requires updating the complement spectra of all vertices connected to $v$. The propagation speed is limited by the graph diameter — the longest shortest path. For the infinite (continuum) limit of ADE chains, this gives a finite maximum propagation speed.

This speed is $c$, a structural consequence of graph consistency under complement-deformation rather than a free parameter.

\subsection*{7.2 Proper time}

For a vertex moving through complement-space with boost parameter $\eta$, the proper time between complement-updates is (M13 exp\_10):

\[d\tau = \frac{dt}{\cosh(\eta)}\]

This is the standard time dilation formula. At $\eta = 0$ (stationary): $d\tau = dt$. As $\eta \to \infty$ (approaching $c$): $d\tau \to 0$.

The graph-theoretic origin: boosted vertices must update more complement relations per unit time, leaving less "budget" for internal dynamics. Time dilation here is complement-budget exhaustion.

\begin{center}\rule{0.5\linewidth}{0.4pt}\end{center}

\section*{8. The complete derivation chain}

The full chain from nothing to proper time, each link verified computationally:

\begin{enumerate}\tightlist
\item \textbf{Self-loop} $\to$ instability under self-reference (M7 exp\_02)
\item \textbf{$\varphi$} from cross-scale consistency (M7 exp\_01)
\item \textbf{PAC} from spectral confinement (M10 exp\_14)
\item \textbf{ADE} from Fibonacci arithmetic closure (M12 exp\_01–02)
\item \textbf{Connection} = addition on ADE root lattice (M12 exp\_01)
\item \textbf{Complement} spectrum uniquely identifies vertices (M13 exp\_01–02)
\item \textbf{Parallax}: different observers, different complements (M13 exp\_02)
\item \textbf{Weyl group} from complement-transformations (M13 exp\_04)
\item \textbf{SEC complexification}: $\mathbb{Z}_2 \to \text{SL}(2,\mathbb{C}) \to \text{SO}(3,1)$ (M12 exp\_10, M13 exp\_05)
\item \textbf{$ds^2$} from Killing form uniqueness (M13 exp\_09)
\item \textbf{$c$} as coherence limit on complement-deformation (M13 exp\_07–08)
\item \textbf{Proper time} $d\tau = dt/\cosh(\eta)$ (M13 exp\_10)
\end{enumerate}

No link in this chain invokes coordinates, manifolds, or metrics. They all emerge from graph combinatorics + SEC dynamics.

\begin{center}\rule{0.5\linewidth}{0.4pt}\end{center}

\section*{9. Honest limitations}

\subsection*{9.1 What works (combined 100/112 = 89\%)}

\begin{itemize}\tightlist
\item M12: 49/52 (94\%) — connection, ADE closure, Lorentz derivation
\item M13 core: 48/52 (92\%) — complement identity, Weyl group, ds², c, proper time
\item M13.5 investigation: 3/8 (38\%) — stress tests revealing structural limits
\end{itemize}

\subsection*{9.2 Failure modes}

\textbf{M12 failures} (3/52):
\begin{enumerate}\tightlist
\item Rate-density proportionality too simplistic — information flow rate is NOT proportional to connection density (exp\_04 T3)
\item Info-Fiedler not proportional — algebraic connectivity doesn't scale linearly with information capacity (exp\_06 T2)
\item Basin depth needs physical coupling — the thermodynamic mechanism works qualitatively but the absolute depth requires a coupling constant not derived from pure graph theory (exp\_07 T4)
\end{enumerate}

\textbf{M13/13.5 failures} (5 documented modes; 9/60 tests failed):
\begin{enumerate}\tightlist
\item Complement-to-Lie PSD: orbit Gram matrix is positive semi-definite but NOT positive definite for all ADE types (exp\_14) — proven fundamental, no invariant metric can fix it. \textbf{Resolved by M14} (orbit Hilbert space makes PSD into a feature)
\item Rapidity composition: discrete graph version has systematic gap from continuous formula (exp\_08 T3)
\item Coherence limit NOT universal: works for A-family, fails for E-family (exp\_15)
\item A-family spectral oscillation at high rank (exp\_17 T3)
\item Random graphs more constrained than ADE at large rank (exp\_17 T4)
\item Alternative complement metrics: no alternative inner product distinguishes same-orbit vertices (exp\_16, 0/4) — the PSD degeneracy is metric-independent, confirming it is fundamental (resolved by the orbit construction of M14)
\end{enumerate}

The framework has two layers: algebraic (solid, \textasciitilde{}92\%) and metric/continuum (weak, \textasciitilde{}38\% under stress). The bridge between them is the ADE classification theorem.

\subsection*{9.3 Structural vs derived}

Approximately 40\% of the M13 results are structural — they pass because complement is defined to satisfy the tested property. The remaining 60\% are derived — they follow from the structure but were not guaranteed a priori. We mark structural results explicitly in the experiment scorecards.

\begin{center}\rule{0.5\linewidth}{0.4pt}\end{center}

\section*{10. Predictions}

\subsection*{10.1 From M12 (8 predictions)}

\begin{longtable}[]{@{}lll@{}}
\toprule
\# & Type & Statement \\
\midrule\endhead
P1 & Precise & SU(2) and SU(3) are the only Fibonacci-compatible ADE gauge groups \\
P2 & Precise & $F_7 = 13 = 1 + 3 + 8 + 1$ gauge dimension closure \\
P3 & Precise & Force relaxation ratios follow $\varphi^{d_2 - d_1}$ \\
P4 & Precise & SEC complexification of $A_1$ gives exactly SO(3,1) \\
P5 & Directional & Laws converge as basin attractors, not imposed rules \\
P6 & Directional & Basin depth correlates with force stability \\
P7 & Constraint & No non-ADE graph type achieves Fibonacci arithmetic closure \\
P8 & Constraint & Complexification of higher ADE types gives larger groups, not Lorentz \\
\bottomrule
\end{longtable}

\subsection*{10.2 From M13 (8 predictions)}

\begin{longtable}[]{@{}lll@{}}
\toprule
\# & Type & Statement \\
\midrule\endhead
P9 & Precise & Complement spectrum uniquely identifies ADE vertices (all types $\leq$ rank 8) \\
P10 & Precise & Complement-transforms form exactly the Weyl group \\
P11 & Precise & Killing form signature = $(-,+,+,+)$ for complexified $A_1$ \\
P12 & Precise & Proper time formula $d\tau = dt/\cosh(\eta)$ from graph boost \\
P13 & Directional & Curvature emerges from connection-density gradients \\
P14 & Directional & Different graph topologies give different "spacetimes" \\
P15 & Constraint & No alternative to Killing form gives Lorentz-invariant bilinear \\
P16 & Constraint & PSD orbit Gram matrix is fundamental (not a defect) \\
\bottomrule
\end{longtable}

\begin{center}\rule{0.5\linewidth}{0.4pt}\end{center}

\section*{11. Relationship to prior papers}

Paper 10 depends on:
\begin{itemize}\tightlist
\item \textbf{Paper 1} (Erasure): PAC conservation as the axiom
\item \textbf{Paper 7} (Symmetry): Self-reference generates $\varphi$ and ADE
\item \textbf{Paper 8} (Quantum Gravity): Response-time framework for forces
\end{itemize}

Paper 10 enables:
\begin{itemize}\tightlist
\item \textbf{Paper 11} (Quantum Mechanics): Orbit Hilbert space on the same ADE graphs
\item Future work on curved spacetime from ADE deformations
\end{itemize}

\begin{center}\rule{0.5\linewidth}{0.4pt}\end{center}

\section*{12. Conclusion}

The Lorentz group is not a postulate; it follows once the simplest possible graph symmetry ($A_1$, two nodes connected by an edge) acquires an entropy arrow (SEC). The metric signature is forced by the Killing form, and the speed of light is the coherence limit of complement-deformation.

The ADE classification — a theorem of pure mathematics, proven in the 1890s — turns out to encode not just the gauge symmetries of particle physics (Paper 7) but also the structure of spacetime itself. Connection corresponds to addition on the root lattice, identity to the complement spectrum, and spacetime geometry to graph symmetry under the SEC arrow.

\textbf{Source code}: \texttt{experiments/milestones/milestone12/} (13 experiments), \texttt{experiments/milestones/milestone13/} (15 experiments)

\begin{center}\rule{0.5\linewidth}{0.4pt}\end{center}

\section*{References}

\begin{enumerate}\tightlist
\item Groom, P. (2026). PACSeries v0.2: Dawn Field Theory. Zenodo. DOI: 10.5281/zenodo.18743674
\item Humphreys, J.E. (1972). Introduction to Lie Algebras and Representation Theory. Springer.
\item Hazewinkel, M. et al. (2010). Algebras, Rings and Modules: Lie Algebras and Hopf Algebras. AMS.
\item McKay, J. (1980). Graphs, singularities, and finite groups. Proc. Symp. Pure Math. 37, 183–186.
\end{enumerate}

\end{document}
